\section{Scattering apparatus and Wigner--Smith time delay}
\label{sec:wigner_smith}

This section states a minimal scattering interface that turns an operational delay measurement into a calibrated proxy for routing overhead.
The goal is to provide a concrete experimental and numerical verification route for lapse ratios without requiring direct access to the underlying compilation log.

\subsection{Scattering apparatus (minimal Hamiltonian model)}
\label{subsec:scattering_apparatus}

Fix an internal (finite) region with Hamiltonian $H$ coupled to $M$ asymptotic channels (leads).
In a standard effective description, the coupling is encoded by a matrix $W$ mapping channel states into the internal region, leading to an effective non-Hermitian Hamiltonian
\begin{equation}
  H_{\mathrm{eff}}(E) := H - \iu\pi\,W W^\dagger,
\end{equation}
and an on-shell scattering matrix of the form \cite{MahauxWeidenmueller1969,Datta1995}
\begin{equation}
  S(E) = \ind - 2\pi\iu\,W^\dagger\left(E - H_{\mathrm{eff}}(E)\right)^{-1}W.
  \label{eq:weidenmueller_S}
\end{equation}
This model covers, at the level of interfaces, a broad class of platforms: mesoscopic transport, microwave networks, photonic structures, and circuit-QED scattering measurements.

\subsection{Wigner--Smith operator and time delay}
\label{subsec:wigner_smith_def}

Consider an $M$-channel scattering apparatus with on-shell scattering matrix $S(E)\in U(M)$.
The Wigner--Smith time-delay operator \cite{Wigner1955,Smith1960} is
\begin{equation}
  Q(E)\ :=\ -\iu\,\hbar\,S(E)^\dagger\,\frac{\dd S}{\dd E}(E),
  \label{eq:wigner_smith_Q}
\end{equation}
and the mean Wigner--Smith delay is
\begin{equation}
  \tau_{\mathrm{WS}}(E)\ :=\ \frac{1}{M}\Tr Q(E).
  \label{eq:wigner_smith_tau}
\end{equation}

\subsection{Resonance linewidth calibration}
\label{subsec:linewidth_calibration}

Near an isolated resonance at energy $E_0$ with linewidth $\gamma$, the scattering matrix admits a Breit--Wigner-type parametrization.
In the single-resonance regime, the peak time delay scales as
\begin{equation}
  \tau_{\mathrm{WS}}(E_0)\ \approx\ \frac{4\hbar}{\gamma},
  \label{eq:ws_linewidth_relation}
\end{equation}
up to apparatus-dependent prefactors that can be calibrated by a reference region.
\par
\noindent For a single channel ($M=1$) one may write $S(E)=\e^{2\iu\delta(E)}$.
In the Breit--Wigner approximation the phase shift takes the standard form
\begin{equation}
  \delta(E)\ \approx\ \arctan\!\left(\frac{\gamma/2}{E_0-E}\right),
\end{equation}
see, e.g., \cite{BreitWigner1936}.
so that $\tau_{\mathrm{WS}}(E)=2\hbar\,\dd\delta/\dd E$ and $\tau_{\mathrm{WS}}(E_0)=4\hbar/\gamma$, consistent with \eqref{eq:ws_linewidth_relation}.

\begin{remark}[Channel conventions]
\label{rem:ws_conventions}
The prefactor in \eqref{eq:ws_linewidth_relation} depends on the convention used for $\tau_{\mathrm{WS}}$ (mean over channels, partial delay, or a particular incoming state) and on how $\gamma$ is defined from the pole structure.
The present paper fixes the convention by \eqref{eq:wigner_smith_Q}--\eqref{eq:wigner_smith_tau} and uses \eqref{eq:ws_linewidth_relation} as a benchmark in the single-resonance regime \cite{BreitWigner1936,Wigner1955,Smith1960}.
\end{remark}

\subsection{From delay to routing overhead}
\label{subsec:ws_to_kappa}

Define the dimensionless delay proxy
\begin{equation}
  \kappa_{\mathrm{WS}}(E)\ :=\ \frac{\tau_{\mathrm{WS}}(E)}{\tau_0},
\end{equation}
where $\tau_0$ is the reference tick duration.
The interface assumption is that a localized resonance probe at location $x$ has a characteristic resonance energy $E_0(x)$ and linewidth $\gamma(x)$ such that
\begin{equation}
  \kappa(x)\ \approx\ \kappa_{\mathrm{WS}}\!\left(E_0(x)\right)
  \approx\ \frac{4\hbar}{\gamma(x)\tau_0}.
  \label{eq:ws_kappa_identification}
\end{equation}
This identification makes lapse ratios directly testable:
\begin{equation}
  \frac{\lapse(x_1)}{\lapse(x_2)}
  = \frac{\kappa(x_2)}{\kappa(x_1)}
  \approx \frac{\gamma(x_1)}{\gamma(x_2)}.
\end{equation}
In practice one calibrates the overall scale by a reference region $x_0$:
define $\gamma_0:=\gamma(x_0)$ and $\kappa_0:=4\hbar/(\gamma_0\tau_0)$ in the same single-resonance convention, so that $\kappa(x)/\kappa_0\approx \gamma_0/\gamma(x)$ within the calibrated bandwidth window.
Appendix~\ref{app:ws_interface_details} records standard trace identities, basis-invariance properties, and loss/dispersion considerations relevant for turning measured $S$-data into calibrated $\kappa$ ratios.
Section~\ref{sec:benchmarks} records consistency targets, and Appendix~\ref{app:adm_details} comments on how probe localization interacts with coarse-graining.


